\section{The Relational--Cyclic Framework}\label{sec:relational_framework}

The SS-PEG framework does not begin with a variational principle. It begins with a simple ontological claim: \textbf{relations are primary, and objects are derivative.} From this claim, the entire architecture of physical theory follows --- not as an assumption, but as a consequence of how relations close into cycles to produce invariants.

This section develops the relational--cyclic framework that underpins the variational principle. The variational principle is the \textit{tool} by which we test whether the framework can generate the gauge symmetries of the Standard Model --- but the framework itself is the \textit{idea}.

\subsection{The Evolution Graph as Fundamental Object}\label{sec:evolution_graph}

The fundamental object in the SS-PEG framework is the \textbf{evolution graph} $\mathcal{G} = (V, E, \mathcal{H}, \{U_e\})$, a quantum graph defined as:

\begin{itemize}
\item \textbf{Vertices} $V$: states of the system (configurations, positions, phase-space points).
\item \textbf{Edges} $E$: relations --- unitary operators $U_e: \mathcal{H}_v \to \mathcal{H}_w$ connecting states.
\item \textbf{Hilbert space} $\mathcal{H} = \bigotimes_{v \in V} \mathcal{H}_v$: the tensor product of local state spaces.
\item \textbf{Edge operators} $\{U_e\}$: parallel transport between states.
\end{itemize}

This graph does \textit{not} represent physical 3D space. It is the abstract space of all possible states of the system --- the ``space of states'' in SS-PEG. Edges represent allowed transitions (evolution operators). The graph is the \textbf{relational substrate}: it encodes what is possible but does not prescribe what is actual.

The evolution graph is the mathematical realization of the ontological claim that relations are primary. Before any ``particle,'' ``field,'' or ``force'' exists, there is only the graph: the web of possible connections between states. Physical entities emerge as \textbf{stable patterns} in this web.

\subsection{The Potential $\to$ Relation $\to$ Cycle $\to$ Invariant Chain}\label{sec:prci_chain}

The SS-PEG framework is built on a single ontological chain:

\begin{equation}
\text{Potential} \;\to\; \text{Relation} \;\to\; \text{Cycle} \;\to\; \text{Invariant}
\end{equation}

\begin{itemize}
\item \textbf{Potential} (P): the primary relational substrate. In gauge theory, this is the vector potential $\mathbf{A}$; in general relativity, the metric tensor $g_{\mu\nu}$; in quantum mechanics, the Lagrangian $L$. The potential defines the \textit{possibility} of relation but does not itself contain observable content.

\item \textbf{Relation} (R): the local prescription for parallel transport. In gauge theory, the covariant derivative $D_\mu = \partial_\mu + A_\mu$; in GR, the Levi-Civita connection $\Gamma$. The relation defines how a state at one point corresponds to a state at an adjacent point.

\item \textbf{Cycle} (C): a closed path through the manifold. In gauge theory, a Wilson loop $W(\gamma) = \mathcal{P}\exp\oint_\gamma A$; in quantum mechanics, an adiabatic cycle in parameter space; in GR, a closed spacetime loop. The cycle is the \textit{mechanism of closure}: only through completing a cycle can the gauge-dependent potential be resolved into a gauge-invariant quantity.

\item \textbf{Invariant} (I): the objective residue of the cycle. In gauge theory, the magnetic flux $\Phi_B$; in quantum mechanics, the Berry phase $\gamma_n$; in GR, the curvature scalar $R$. The invariant is the only observer-independent reality --- the part of the potential that persists across all coordinate systems and gauge choices.
\end{itemize}

This chain is not a metaphor. It is a precise mathematical isomorphism that appears in every pillar of modern physics. The following sections demonstrate this through four foundational phenomena.

\subsection{The Aharonov--Bohm Effect: Potential as Primary Reality}\label{sec:aharonov_bohm}

The Aharonov--Bohm (AB) effect stands as the quintessential example of the physicality of potentials. In the standard experiment, a charged particle is split into two paths passing on opposite sides of a long solenoid. The magnetic field $\mathbf{B}$ is confined entirely within the solenoid, meaning the particle moves through a region where $\mathbf{B} = 0$. Despite the absence of local magnetic forces, the particle acquires a phase shift depending on the magnetic flux $\Phi_B$ enclosed by its paths.

In the relational view, the vector potential $\mathbf{A}$ is the fundamental potential that defines the local phase relations of the wave function. The ``force-free'' region is not empty; it is saturated with relational information encoded in $\mathbf{A}$. The phase relation between two points $a$ and $b$ along a path $\gamma$ is:

\begin{equation}
\Delta\phi = \frac{q}{\hbar} \int_\gamma \mathbf{A} \cdot d\mathbf{l}
\end{equation}

The AB effect is only observable when the two paths recombine, forming a closed cycle $\mathcal{C}$. The interference pattern at the detector is the physical manifestation of the holonomy of the connection around that cycle:

\begin{equation}
\oint_\mathcal{C} \mathbf{A} \cdot d\mathbf{l} = \iint_S (\nabla \times \mathbf{A}) \cdot d\mathbf{S} = \Phi_B
\end{equation}

The cycle acts as a topological probe. Even though the particle never enters the region where the curvature (the magnetic field) is non-zero, the cycle ``winds'' around the singularity, capturing the invariant magnetic flux. This illustrates a fundamental principle: \textbf{the invariant does not require local interaction with a force; it only requires the completion of a cycle that encloses the source of the potential's non-triviality.}

\subsection{Berry Phase: Geometric Memory of Cycles}\label{sec:berry_phase}

The Berry phase extends the logic of the Aharonov--Bohm effect from the movement of particles in physical space to the movement of states in parameter space. When a quantum system's Hamiltonian $H(\mathbf{R})$ depends on parameters $\mathbf{R}$ varied slowly (adiabatically) along a cycle $\mathcal{C}$, the wave function acquires a phase factor that is purely geometric:

\begin{equation}
\gamma_n(\mathcal{C}) = \oint_\mathcal{C} \mathbf{A}_n \cdot d\mathbf{R} = \iint_S \mathbf{F}_n \cdot d\mathbf{S}
\end{equation}

where $\mathbf{A}_n(\mathbf{R}) = i\langle n(\mathbf{R}) | \nabla_\mathbf{R} | n(\mathbf{R}) \rangle$ is the Berry connection and $\mathbf{F}_n = \nabla \times \mathbf{A}_n$ is the Berry curvature.

The Berry phase is an invariant that appears only when parameters return to their initial configuration, completing a cycle. During this cyclic evolution, the wave function acquires a ``memory'' of the path taken. This memory is not stored in the Hamiltonian or the energy levels (which would yield a dynamical phase) but in the curvature of the quantum bundle itself.

For non-cyclical variations, the Berry phase can be made to vanish by re-phasing the eigenstates at each point (a gauge transformation). For a cycle, the phase becomes an invariant property that cannot be cancelled. The integral of the Berry curvature over a closed surface yields a \textbf{Chern number} --- an integer topological invariant classifying the bundle structure. These invariants underlie quantized phenomena such as the Quantum Hall Effect.

\subsection{Lagrangian Mechanics: Action as Relational Constraint}\label{sec:lagrangian}

Lagrangian mechanics provides a unified way to derive equations of motion for any system. The configuration space $M$ is the manifold of all possible positions $q$. The tangent bundle $TM$ adds velocities $\dot{q}$. The Lagrangian $L(q, \dot{q}, t)$ is a potential function summarizing the entire dynamics of the system.

The physical path is not determined by forces but by a relational requirement: the action $S = \int L\, dt$ must be stationary. The Euler--Lagrange equations:

\begin{equation}
\frac{d}{dt}\left(\frac{\partial L}{\partial \dot{q}_j}\right) - \frac{\partial L}{\partial q_j} = 0
\end{equation}

are the local expressions of this relational balance. They define how the state of the system at time $t$ must relate to its state at $t+dt$ to ensure that the global action remains invariant under small perturbations.

The Lagrangian is not unique: $L \to L + dF/dt$ leaves the equations invariant. This is a \textbf{gauge transformation of the Lagrangian potential}. The ``relation''---the dynamics---is the invariant part of this potential that survives such shifts.

Lagrangian mechanics also excels at handling constraints. \textbf{Holonomic constraints} connect coordinates and can be expressed as $g(q, t) = 0$, restricting the system to a sub-manifold. \textbf{Non-holonomic constraints}, which depend on velocities and cannot be integrated, introduce path-dependence. In such systems, completing a cycle in one set of variables can lead to a net change in another set --- the classical manifestation of holonomy. A rolling sphere that completes a cycle on a flat surface may return to its starting point with a different orientation. Here, the ``potential'' is the set of constraint relations, the ``cycle'' is the path on the surface, and the ``invariant'' is the resulting rotation.

\subsection{General Relativity: Metric as Ultimate Relational Potential}\label{sec:gr}

In Einstein's theory, the ``force'' of gravity is entirely replaced by the curvature of spacetime. The central object is the \textbf{metric tensor} $g_{\mu\nu}$, the gravitational potential. It defines the interval $ds^2 = g_{\mu\nu}\, dx^\mu\, dx^\nu$, the ultimate relational measure of distance and time between events.

While the metric is the potential, the \textbf{Levi-Civita connection} $\nabla$ is the system of relations. The connection defines parallel transport: it tells us how a vector at one point in spacetime relates to a vector at a neighboring point. A vector $V^\mu$ is parallel-transported along a curve $\gamma$ if:

\begin{equation}
\frac{dV^\mu}{d\tau} + \Gamma^\mu_{\alpha\beta} V^\alpha \frac{dx^\beta}{d\tau} = 0
\end{equation}

The Christoffel symbols $\Gamma$ act as the ``connection coefficients'' mediating these local relations.

In General Relativity, gravity manifests as the \textbf{non-closure of relations around a cycle}. When a vector is parallel-transported around an infinitesimal closed loop, the discrepancy between initial and final vectors is determined by the \textbf{Riemann curvature tensor}:

\begin{equation}
R^\rho{}_{\sigma\mu\nu} = \partial_\mu \Gamma^\rho_{\nu\sigma} - \partial_\nu \Gamma^\rho_{\mu\sigma} + \Gamma^\rho_{\mu\lambda}\Gamma^\lambda_{\nu\sigma} - \Gamma^\rho_{\nu\lambda}\Gamma^\lambda_{\mu\sigma}
\end{equation}

The curvature is the ``field strength'' of the gravitational potential, just as $\mathbf{B}$ is the field strength of $\mathbf{A}$. The holonomy of the Levi-Civita connection around a cycle $\mathcal{C}$ measures the net rotation and scaling of vectors. For flat spacetime (Minkowski metric), the holonomy group is trivial; for curved spacetime, the holonomy group reflects the independent degrees of freedom of the gravitational field.

The metric components $g_{\mu\nu}$ depend on the choice of local coordinates (the observer's ``gauge''), but the \textbf{curvature invariants} constructed from the metric --- the Ricci scalar $R$, the Weyl tensor components --- do not. These invariants represent the absolute, coordinate-independent reality of the gravitational field.

\subsection{Structural Unity Across Domains}\label{sec:structural_unity}

The four pillars of modern physics, when viewed through the relational lens, reveal a profound structural identity:

\begin{equation}
\begin{array}{lcccc}
\hline
\textbf{Feature} & \textbf{AB Effect} & \textbf{Berry Phase} & \textbf{Lagrangian Mech.} & \textbf{General Relativity} \\
\hline
\text{Potential} & A_\mu & A_n & L & g_{\mu\nu} \\
\text{Relation} & \text{Phase transport} & \text{Adiabatic transition} & \text{Variational stability} & \text{Affine connection} \\
\text{Cycle} & \text{Spatial loop} & \text{Parameter loop} & \text{Periodic orbit} & \text{Spacetime loop} \\
\text{Invariant} & \Delta\phi & \gamma_n & \text{Conserved qty} & R \\
\text{Curvature} & \mathbf{B} & \mathbf{F}_n & \text{E--L eqs.} & R^\rho{}_{\sigma\mu\nu} \\
\hline
\end{array}
\end{equation}

Every column is a \textbf{different projection} of the same relational graph $\mathcal{G}$ into a different physical context. The potentials are the primary objects; relations connect them locally; cycles probe global structure; and invariants are the only observer-independent reality.

\begin{quote}
\textbf{``The cycle is all you need.''} From coupling constants (master formula $\alpha_X = e^{S_X}/4\pi$ with direction vectors read from cycle structure) to spacetime curvature (Riemann $=$ cycle discrepancy), everything physical is an invariant of a cycle in the relational graph.
\end{quote}

\subsection{Relational Density as Order Parameter}\label{sec:relational_density}

The evolution graph $\mathcal{G}$ can be dense (many connections) or sparse (few). The \textbf{relational density} $\rho(v)$ of a vertex $v$ measures how many independent cycles pass through $v$ relative to the total number of possible connections:

\begin{equation}
\rho(v) = \frac{\#\{\text{cycles of length $\leq L$ through } v\}}{\text{Vol}(B_L(v))}
\end{equation}

where $B_L(v)$ is the ball of radius $L$ centered at $v$ (vertices reachable in $\leq L$ steps).

In entangled quantum systems, this density is equivalent to the \textbf{entanglement entropy}:

\begin{equation}
\rho(v) = S_{\mathrm{ent}}(\rho_v) = -\mathrm{Tr}(\rho_v \log \rho_v)
\end{equation}

where $\rho_v$ is the reduced density matrix of the total state restricted to vertex $v$.

The relational density is the \textbf{order parameter} that controls which gauge symmetries emerge in different regions of the graph. Low density ($\rho < \rho_1$): few independent cycles, holonomies commute, U(1) emerges. Medium density ($\rho_1 < \rho < \rho_2$): sufficient cycles for non-commutativity, su(2) emerges. High density ($\rho > \rho_2$): dense network of interlocked cycles, su(3) emerges.

This hierarchy is not imposed --- it emerges from the topology of the graph. The phase transition between these regimes is sharp: there is no ``partial SU(3)'' at intermediate energies. The gauge group is either exact or absent. This is the physical meaning of the density phase transition demonstrated in Section~\ref{sec:density}.

\subsection{The Cartan--Weyl Graph and the Variational Principle}\label{sec:cw_graph}

The Cartan--Weyl (CW) root graph provides the bridge between the continuous algebraic structure (Killing metric) and the discrete graph topology (relational density). For a simple Lie algebra $\mathfrak{g}$, the CW graph $G_{\mathrm{CW}} = (V, E)$ has:

\begin{itemize}
\item Vertices $V = \Phi$: one for each root.
\item Edges $(\alpha, \beta) \in E$ iff $\langle \alpha, \beta^\vee \rangle \neq 0$, where $\beta^\vee = 2\beta / (\beta, \beta)$ is the coroot.
\end{itemize}

The CW graph encodes the root system geometry: edges connect roots that are non-orthogonal, and edge weights are the Cartan integers. For the Standard Model gauge algebra $\mathfrak{su}(3) \oplus \mathfrak{su}(2) \oplus \mathfrak{u}(1)$, the CW graph has two disconnected components corresponding to the SU(3) and SU(2) factors.

The spectrum of the adjacency matrix $A$ of $G_{\mathrm{CW}}$ determines the \textbf{Ramanujan ratio} $\mathcal{R}$, a measure of the spectral expansion quality of the graph. A graph is Ramanujan if all non-trivial eigenvalues $\lambda$ satisfy $|\lambda| \leq 2\sqrt{d_{\mathrm{avg}} - 1}$, where $d_{\mathrm{avg}}$ is the average degree. The Ramanujan property is equivalent to optimal spectral expansion and maximal chaos in the corresponding quantum system.

\begin{theorem}[Ramanujan Property of SM Factors]
The CW graphs of $\su(2)$ and $\su(3)$ are Ramanujan with ratios $\mathcal{R}(\su(2)) = 1/2$ and $\mathcal{R}(\su(3)) = (\sqrt{57} - 3)/(2\sqrt{17}) \approx 0.5523$, respectively. Both pass all six alternative definitions of the Ramanujan property unanimously. The product $\su(3) \oplus \su(2)$ is also Ramanujan (ratio 0.657).
\end{theorem}

The Ramanujan property is not a mathematical curiosity. It is a \textbf{necessary condition for spectral optimality}: the SM factors are deeply Ramanujan because the variational principle selects graphs with optimal spectral expansion. Paper II develops this connection in detail.

\subsection{From Relational Ontology to Variational Principle}\label{sec:bridge_to_variational}

The relational--cyclic framework establishes the \textbf{ontological foundation} of the SS-PEG program: reality is not made of things but of relations. The evolution graph is the mathematical realization of this ontology.

The \textbf{Killing metric variational principle} is the \textbf{test} of this ontology: it asks whether the relational structure encoded in the evolution graph can generate the gauge symmetries of the Standard Model. The Killing metric $K_{ab} = \mathrm{Tr}(\mathrm{ad}_a \circ \mathrm{ad}_b)$ is the \textit{meta-potential} of the framework: it measures the intensity of mutual interference between all pairs of generators, encoding the global geometric pattern of how all commutators relate.

The variational principle is not the \textit{idea} --- it is the \textit{tool}. The idea is the relational--cyclic ontology. The tool is the Killing metric minimization. Together, they form the complete SS-PEG framework:

\begin{equation}
\underbrace{\text{Evolution graph} \to \text{Cycles} \to \text{Killing metric}}_{\text{Ontology (idea)}} \;\to\; \underbrace{\text{Minimization of } T_K}_{\text{Variational principle (tool)}} \;\to\; \underbrace{\text{Lie algebra emergence}}_{\text{Prediction}}
\end{equation}
